\documentclass[12pt,a4paper,UTF8]{ctexart}
\usepackage[a4paper,top=22mm,bottom=22mm,left=22mm,right=22mm,headheight=15pt]{geometry}
\usepackage{amsmath,amssymb,mathtools,bm}
\usepackage{graphicx,xcolor,booktabs,longtable,array,enumitem,fancyhdr,caption,float,microtype,setspace}
\newcommand{\Needspace}[1]{\par\begingroup\dimen0=#1\relax\dimen2=\pagegoal\advance\dimen2 by-\pagetotal\ifdim\dimen0>\dimen2\newpage\fi\endgroup}
\newcommand{\needspace}[1]{\par\begingroup\skip0=\lastskip\vskip-\skip0\dimen0=#1\relax\advance\dimen0 by\skip0\vskip\dimen0\penalty-100\vskip-\dimen0\vskip\skip0\penalty9999\endgroup}
\usepackage[hidelinks,unicode]{hyperref}
\usepackage[most]{tcolorbox}
\definecolor{ContestBlue}{HTML}{17365D}
\hypersetup{pdftitle={第 044 套},pdfauthor={}}
\setstretch{1.24}
\setlength{\parindent}{2em}
\setlength{\parskip}{0.35em}
\setlength{\emergencystretch}{2em}
\setlist{itemsep=0.35em,topsep=0.4em}
\captionsetup{font=small,labelfont=bf,justification=centering}
\setcounter{secnumdepth}{0}
\pagestyle{fancy}
\fancyhf{}
\fancyhead[L]{\small 第 044 套}
\fancyfoot[C]{\thepage}
\renewcommand{\headrulewidth}{0.4pt}
\ctexset{
  section={format=\Large\bfseries\color{ContestBlue},beforeskip=1.2em,afterskip=0.7em},
  subsection={format=\large\bfseries\color{ContestBlue},beforeskip=1.0em,afterskip=0.5em},
  subsubsection={format=\normalsize\bfseries,beforeskip=0.8em,afterskip=0.35em}
}
\makeatletter
\newcommand{\tightlist}{\setlength{\itemsep}{0pt}\setlength{\parskip}{0pt}}
\newcommand{\maxwidth}{\ifdim\Gin@nat@width>\linewidth\linewidth\else\Gin@nat@width\fi}
\newcommand{\maxheight}{\ifdim\Gin@nat@height>0.72\textheight 0.72\textheight\else\Gin@nat@height\fi}
\newcommand{\pandocbounded}[1]{\begingroup\setkeys{Gin}{width=\maxwidth,height=\maxheight,keepaspectratio}#1\endgroup}
\makeatother
\begin{document}
\begin{titlepage}
  \centering
  \vspace*{58mm}
  {\Huge\bfseries\color{ContestBlue} 第 044 套\par}
  \vfill
  \vspace*{28mm}
\end{titlepage}
\section{第 1 题}\label{ux7b2c-1-ux9898}

足够大的斜面倾角为 \(\theta\)，木块与斜面之间的动摩擦系数为
\(\mu=\tan\theta\)。在斜面内沿水平等高线方向给木块初速度
\(v_0\)。求木块最终的速度 \(v'\)。取 \(0<\theta<\pi/2\)，忽略空气阻力。

\begin{figure}[H]
\centering
\includegraphics[width=0.86\linewidth,height=\textheight,keepaspectratio,alt={题1　题设与符号}]{../figures/question-01.pdf}
\caption{题1　题设与符号}
\end{figure}

\newpage

\section{第 2 题}\label{ux7b2c-2-ux9898}

质量 \(m\) 的粒子以初速大小 \(v_0\) 在光滑平面轨道 \(r=r_0e^{-a\theta}\)
上运动，\(a>0\)，初始
\(t=0,\theta=0\)。无其他做功的力。求：（1）速度与向外径向单位向量的夹角
\(\beta\)；（2）关于原点的角动量变化率；（3）轨道作用力。分别注明沿
\(\theta\) 增大和减小方向的情况。

\begin{figure}[H]
\centering
\includegraphics[width=0.86\linewidth,height=\textheight,keepaspectratio,alt={题2　题设与符号}]{../figures/question-02.pdf}
\caption{题2　题设与符号}
\end{figure}

\newpage

\section{第 3 题}\label{ux7b2c-3-ux9898}

假想磁单极子的磁荷为 \(q\)，受力 \(F=qB\)，磁场
\(B=\mu_0q/(4\pi r^2)\)。它在竖直平面内沿开口向上的光滑U形轨道运动，底部穿过一匝圆线圈中心，局部轨道沿线圈轴线。线圈半径
\(b\) 远小于U形轨道宽度，导线半径 \(a\ll b\)，电阻率为
\(\rho\)。单极子质量为 \(m\)，从高于底部 \(H\) 处静止释放，忽略自感。

求：（1）单极子距线圈中心轴向坐标为 \(x\) 时的磁通；（2）轴向速度为
\(v\)
时的电动势；（3）一次通过的速度变化；（4）通过多少次才能静止；（5）能量法与动力学法是否一致。可用
\(\int_{-\infty}^{\infty}(1+u^2)^{-3}du=3\pi/8\)。

\begin{figure}[H]
\centering
\includegraphics[width=0.86\linewidth,height=\textheight,keepaspectratio,alt={题3　题设与符号}]{../figures/question-03.pdf}
\caption{题3　题设与符号}
\end{figure}

\newpage

\section{第 4 题}\label{ux7b2c-4-ux9898}

密闭圆柱容器固定在可沿水平面自由移动的车上，空容器、车板、车轮总质量
\(M=150\,\mathrm{kg}\)，容器长 \(L=10\,\mathrm m\)，内有质量
\(m=120\,\mathrm{kg}\) 的理想气体。初温 \(T_0=300\,\mathrm K\)，压强
\(P_0=150\,\mathrm{atm}\)。左端保持 \(300\,\mathrm K\)，右端加热并保持
\(332\,\mathrm K\)，其他表面绝热。稳态温度沿长度近似线性，无化学反应。求稳态压强和容器自身质心的位移。忽略外部水平作用力及热流携带的净水平动量。

\begin{figure}[H]
\centering
\includegraphics[width=0.821\linewidth,height=\textheight,keepaspectratio,alt={题4　题设与符号}]{../figures/question-04.pdf}
\caption{题4　题设与符号}
\end{figure}

\newpage

\section{第 5 题}\label{ux7b2c-5-ux9898}

柱状物体的横截面表面是非凹镜面，可含平面、凸面和棱角。物体绕内部轴旋转，固定激光束通过转轴入射。记录转角
\(\theta\) 与有向出射偏角
\(\alpha\)，原图可整理为以下开区间上的线段（角度单位为度）： \[
\alpha=
\begin{cases}
2\theta,&0<\theta<45,\\
\theta-90,&45<\theta<135,\\
2(\theta-180),&135<\theta<225,\\
\theta-270,&225<\theta<315,\\
2(\theta-360),&315<\theta<360.
\end{cases}
\]
在45、135、225、315度处有跳变。求：（1）横截面形状；（2）若表面距转轴最远处为
\(r_0\)，写表面的分段极坐标方程。取 \(\theta\) 为逆时针转角，\(\alpha\)
相对入射光反向延长的零度方向计量。

\begin{figure}[H]
\centering
\includegraphics[width=0.815\linewidth,height=\textheight,keepaspectratio,alt={题5　题设与符号}]{../figures/question-05.pdf}
\caption{题5　题设与符号}
\end{figure}

\newpage

\section{第 6 题}\label{ux7b2c-6-ux9898}

质量 \(m=3.18\,\mathrm{kg}\) 的均匀圆柱放在质量
\(M=6.4\,\mathrm{kg}\)、倾角 \(\theta=40.7^\circ\)
的斜面上，\(g=9.8\,\mathrm{m/s^2}\)。分别求：

（1）斜面固定、接触光滑时圆柱加速度 \(a_1\)；
（2）斜面可在光滑水平地面移动、接触光滑时圆柱相对斜面加速度 \(a_2\)；
（3）斜面固定、圆柱纯滚动的最小摩擦系数 \(\mu_{1m}\)；
（4）该纯滚动加速度 \(a_3\)；
（5）斜面可自由移动、圆柱纯滚动的最小摩擦系数 \(\mu_{2m}\)；
（6）第（5）种情况圆柱的加速度
\(a_4\)。最后一问分别给出相对斜面及实验室中的加速度，避免参考系歧义。

\begin{figure}[H]
\centering
\includegraphics[width=0.702\linewidth,height=\textheight,keepaspectratio,alt={题6　题设与符号}]{../figures/question-06.pdf}
\caption{题6　题设与符号}
\end{figure}

\newpage

\section{第 7 题}\label{ux7b2c-7-ux9898}

反应 \(\gamma+p\to K^++\Lambda^0\) 中，
\(M_p=938\,\mathrm{MeV}/c^2\)、\(M_K=494\,\mathrm{MeV}/c^2\)、\(M_\Lambda=1116\,\mathrm{MeV}/c^2\)。
（1）质子静止时求光子阈能； （2）若靶质子有不超过
\(p_F=200\,\mathrm{MeV}/c\)
的费米动量，采用自由在壳质子近似，求最低阈能； （3）\(\Lambda^0\) 以
\(\beta c\) 飞行并衰变成
\(p+\pi^-\)，\(\beta=0.773\)、\(M_\pi=140\,\mathrm{MeV}/c^2\)，求
\(\pi^-\) 最大实验室动量； （4）求 \(\pi^-\) 相对 \(\Lambda^0\)
飞行方向的最大夹角。核结合能和旁观核反冲在第（2）问未给出，另注明这一近似。

\begin{figure}[H]
\centering
\includegraphics[width=0.86\linewidth,height=\textheight,keepaspectratio,alt={题7　题设与符号}]{../figures/question-07.pdf}
\caption{题7　题设与符号}
\end{figure}

\newpage

\section{第 8 题}\label{ux7b2c-8-ux9898}

三同心金属球壳半径为 \(r_1>r_2>r_3\)。外壳电势为
\(U\)，内壳经表面绝缘的导线接地，中壳悬浮。外---中层介质参数为
\(\varepsilon_{r1},\sigma_1\)，中---内层为
\(\varepsilon_{r2},\sigma_2\)。足够长时间后达到电学稳态。求：（1）内壳电流大小；（2）外壳上的自由电荷与其邻接介质的极化电荷总量
\(q_1\)；（3）中壳自由电荷 \(q_{f2}\)。

参数：\(\sigma_1=5.08\,\mathrm{S/m},\sigma_2=4.42\,\mathrm{S/m}\)，
\(\varepsilon_{r1}=2.98,\varepsilon_{r2}=1.16\)，
\(r_1=0.724\,\mathrm m,r_2=0.491\,\mathrm m,r_3=0.216\,\mathrm m\)，
\(\varepsilon_0=8.85\times10^{-12}\,\mathrm{F/m},U=6.19\,\mathrm V\)。
先给出外壳内表面与邻接极化层的电荷；若``外壳总量''包含外表面，另说明外部边界条件。

\begin{figure}[H]
\centering
\includegraphics[width=0.733\linewidth,height=\textheight,keepaspectratio,alt={题8　题设与符号}]{../figures/question-08.pdf}
\caption{题8　题设与符号}
\end{figure}
\end{document}
